% ADL Lite: An Event-First Operational Ontology for Multi-Agent Concept Consensus
% Generated: 2026-05-31 | Synced with MD v27-refs
% Target: Semantic Web Journal (IOS Press)
%
% Contents: \S1-8, 27 references, CRDT proofs, distributed deployment, Nanopub comparison

\documentclass[11pt,a4paper]{article}
\usepackage[utf8]{inputenc}
\usepackage[T1]{fontenc}
\usepackage{amsmath,amssymb,amsthm}
\usepackage{graphicx}
\usepackage{hyperref}
\usepackage{booktabs}
\usepackage{listings}
\usepackage{xcolor}
\usepackage{enumitem}
\usepackage[margin=2.5cm]{geometry}
\usepackage{fancyhdr}
\usepackage{cite}

\newtheorem{definition}{Definition}[section]
\theoremstyle{remark}
\newtheorem{remark}{Remark}[section]

\lstset{basicstyle=\ttfamily\small,breaklines=true,frame=single,xleftmargin=15pt,
  numbers=left,numberstyle=\tiny,language=Python,showstringspaces=false,
  keywordstyle=\color{blue}\bfseries,commentstyle=\color{gray}\itshape,
  stringstyle=\color{red!70!black}}

\pagestyle{fancy}\fancyhf{}
\fancyhead[L]{\small A. Zhou / ADL Lite}
\fancyhead[R]{\small Event-First Operational Ontology}
\renewcommand{\headrulewidth}{0.4pt}

\title{\textbf{ADL Lite: An Event-First Operational Ontology\\for Multi-Agent Concept Consensus\\and Knowledge Graph Authoring}}
\author{Allen Zhou\textsuperscript{1,*}\\\small\textsuperscript{1}CEIEC AI Infrastructure\\\small\textsuperscript{*}Corresponding author}
\date{}

\begin{document}
\maketitle

\begin{abstract}
The construction of knowledge graphs (KGs) increasingly relies on large language model (LLM) agents for discovery and authoring. However, existing ontology frameworks assume object-first design---concepts are mutable entities with stored properties---which creates fundamental tension with agentic workflows: multiple agents produce competing interpretations, status and confidence become stale fields, and centralized trust is required for validation.

We present ADL Lite, an event-first operational ontology that redefines concepts as append-only chains of typed events. Drawing on Wittgenstein's \textit{Tractatus} \S1.1---``the world is the totality of facts, not of things''---ADL Lite treats every knowledge update as a new event rather than a mutation of existing state. The system provides a four-layer, Markdown-native syntax (L1: YAML, L2: Prose, L3: Structured Semantic Anchoring blocks, L4: Typed Action blocks). We introduce three core mechanisms: (1) Structured Semantic Anchoring (SSA), (2) a cryptographic EventChain with chained SHA-256 hashing for tamper evidence, and (3) a multi-agent consensus protocol with fork management and CRDT-inspired offline-first synchronization.

We evaluate ADL Lite through six controlled experiments plus a 10--100 agent distributed deployment simulation, validated by \textbf{350 automated tests} at 81\% code coverage. The framework is implemented in Python 3.10+ with Pydantic v2, SQLite-backed hybrid memory, and a Feishu/Lark bridge. We provide a formal CRDT lattice proof for derived states, an RDF/OWL/PROV-O mapping with 100\% information preservation, and comprehensive comparisons with Nanopublications, RDF-star, and Wikidata's history model.

\vspace{3mm}\noindent\textbf{Keywords:} operational ontology, multi-agent consensus, knowledge graph authoring, event-first design, structured semantic anchoring, cryptographic chain integrity, CRDT
\end{abstract}

% ===================================================================
\section{Introduction}
\label{sec:intro}

Knowledge graphs are increasingly authored not by individual experts but by fleets of large language model (LLM) agents operating in parallel~\cite{pan2024unifying,chen2024agentpoison}. This shift raises a fundamental architectural question: \textit{what is the right primitive for representing knowledge when the authors are autonomous agents?}

Existing ontology frameworks---RDF/OWL~\cite{mcguinness2004owl}, property graphs~\cite{angles2008survey}, Pydantic schemas~\cite{colvin2021pydantic}---use object-first design. Concepts are mutable objects. This breaks down in multi-agent settings:

\begin{enumerate}[label=(\arabic*)]
    \item \textbf{State conflict}: Simultaneous writes produce race conditions.
    \item \textbf{Auditability loss}: Reasons for state changes are lost.
    \item \textbf{Trust centralization}: A coordinator must arbitrate writes.
    \item \textbf{Training burden}: Agents need formal ontology languages.
\end{enumerate}

We argue for \textbf{event-first} design, grounded in Wittgenstein~\cite{wittgenstein1922tractatus}: ``The world is the totality of facts, not of things.'' A concept \textit{is} its event chain. Status, confidence, validators are \textit{computed}, not stored. This eliminates write-time race conditions; it reduces merge-time non-determinism to consistent ordering within timestamp buckets (\S\ref{sec:failure}).

ADL Lite provides: (1) a four-layer Markdown-native syntax, (2) Structured Semantic Anchoring (SSA), (3) a cryptographic EventChain with chained hash linking, (4) a multi-agent consensus protocol with fork management, and (5) a reference implementation with 350 automated tests.

Research questions:
\begin{itemize}
    \item \textbf{RQ1 (Concurrency):} Can event-first design eliminate race conditions?
    \item \textbf{RQ2 (Ambiguity):} Does SSA reduce referential ambiguity?
    \item \textbf{RQ3 (Consensus):} Is CRDT-inspired merge sufficient for offline-first consensus?
\end{itemize}

% ===================================================================
\section{Related Work}
\label{sec:related}

\subsection{Operational Ontology and Event-First Design}

Operational ontology has roots in process philosophy~\cite{whitehead1929process} and event-driven architectures~\cite{fowler2005event}. PROV-O~\cite{lebo2013prov} models provenance as metadata \textit{about} concepts. The Event Ontology~\cite{raimond2007event} and SEM~\cite{vanhage2011sem} model events as first-class but do not replace the concept-as-object paradigm.

\subsection{KG Authoring by LLM Agents}

Pan et al.~\cite{pan2024unifying}, Trajanoska et al.~\cite{trajanoska2023enhancing}, and Chen et al.~\cite{chen2024exploring} used LLMs for KG construction but relied on centralized orchestrators. ADL Lite targets the \textit{authoring} phase---discovery and consensus---not completion or QA (FB15k, WN18RR, OAEI).

\subsection{Consensus in Multi-Agent Systems}

Consensus protocols for distributed systems~\cite{lamport1982byzantine,ongaro2014raft} and multi-agent systems~\cite{wooldridge2009multiagent,jurca2003reputation,brandt2016handbook} assume agents vote on propositions. ADL Lite's model is inspired by CRDT principles~\cite{shapiro2011crdt}---append-only events are trivially mergeable---but derived states use last-writer-wins, not full lattice joins (\S\ref{sec:limitations}).

\subsection{Cryptographic Integrity in Knowledge Systems}

Blockchain~\cite{third2017linkchains} and content-addressable storage~\cite{benet2014ipfs} provide integrity but with overhead. ADL Lite uses lightweight SHA-256 hash chaining per concept.

\subsection{Event-Centric RDF and Provenance Models}

\textbf{Nanopublications}~\cite{groth2010nanopub} model assertions as immutable, signed RDF graphs with provenance. ADL Lite differs by linking events into causal chains with predecessor hash binding, and by using Markdown-native authoring for LLM accessibility.

\textbf{RDF-star}~\cite{hartig2024rdfstar} enables provenance annotations within RDF but does not prescribe lifecycle state machines or fork management.

\textbf{Wikidata}~\cite{vrandecic2014wikidata} records every edit as a revision---event-first thinking at scale. ADL Lite adds cryptographic chain integrity, structured fork resolution, and a portable authoring format.

\textbf{RDF change sets}~\cite{bernerslee2004delta} represent graph deltas. ADL Lite events carry semantic types (REGISTER, VALIDATE, FORK) beyond add/delete, embedding intent and reasoning.

% ===================================================================
\section{The ADL Lite Framework}
\label{sec:framework}

\subsection{Design Philosophy}

\textbf{P1: Event-first.} A concept \textit{is} its \texttt{EventChain}. Properties are computed at query time.

\textbf{P2: Agent-accessible syntax.} Markdown is LLM-trained, human-readable, and supports layered structure.

\textbf{P3: Concurrency-safe.} Append-only + thread-safe locking + CRDT-inspired merge.

\subsection{Four-Layer Syntax}

\textbf{L1: YAML Front Matter.} Identity, type, scope. The L1 is a \textit{snapshot} derived from the EventChain.

\begin{lstlisting}[caption={L1 YAML example}]
adl_type: discovery
adl_id: disc-capital-trap
status: provisional
confidence: 0.84
domain: financial_aml
scope: public
\end{lstlisting}

\textbf{L2: Markdown Body.} Free-form prose with \texttt{[[Wiki Links]]}.

\textbf{L3: Structured Semantic Anchoring.} Typed blocks (\texttt{adl:relation}, \texttt{adl:evidence}, \texttt{adl:seal}) with machine-parseable YAML.

\textbf{L4: Action Blocks.} Typed actions with precondition-enforced side effects, validated against the YAML ontology registry.

\subsection{EventChain: The Core Data Structure}

\begin{definition}[EventChain]
An EventChain $C = (e_0, \ldots, e_n)$ where $e_i = (\mathrm{id}_i, \mathrm{cid}, \tau_i, a_i, r_i, t_i, p_i, \mathrm{prev}_i, h_i)$.
\end{definition}

$\mathrm{id}_i$ is a full 128-bit UUID4 (32 hex characters, $2^{122}$ random bits). $h_i = \mathrm{SHA\text{-}256}(\sigma(e_i) \parallel h_{i-1})$ for $i>0$, $h_0 = \mathrm{SHA\text{-}256}(\sigma(e_0))$, where $\sigma$ serializes to canonical JSON sorted by key. The integrity invariant is $\forall i>0: \mathrm{prev}_i = \mathrm{id}_{i-1} \land h_i = \mathrm{SHA\text{-}256}(\sigma(e_i) \parallel h_{i-1})$. Status derivation: $\mathrm{status}(C) = f(\tau_k)$ where $k = \max\{i \mid \tau_i \in \mathcal{L}\}$ and $\mathcal{L}$ is the lifecycle event set.

\subsection{Structured Semantic Anchoring (SSA)}

SSA uses: (1) \textbf{Pronoun prohibition}---pattern-based detection, 29-sentence benchmark achieving \textbf{precision 0.812, recall 1.000, F1 0.897}; (2) \textbf{Typed slots} via Pydantic models; (3) \textbf{Ontology registry} with 10 predicates, 9 actions, 4 scope prefixes.

% ===================================================================
\section{Multi-Agent Consensus Protocol}
\label{sec:consensus}

\subsection{Lifecycle State Machine}

provisional $\to$ validated/forked $\to$ deprecated $\to$ archived. Valid transitions enforced by the ontology registry.

\subsection{Fork Management}

Jaccard similarity on (relation, target) pairs: $\mathrm{sim}(C_a, C_b) = |E_a \cap E_b| / |E_a \cup E_b|$. MERGE ($\geq 0.90$), PARALLEL ($<0.90$), PRUNE ($>180$ days idle). Threshold sensitivity: zero false positives across 0.70--0.95 sweep, F1=0.667 at default 0.90.

\subsection{Offline-First Edge Synchronization}

CRDT-inspired eventual consistency via \texttt{SyncManager.merge()}: collect, deduplicate by event\_id, sort by (timestamp, hash), rebuild. \texttt{EdgeNode} buffers side effects when offline, drains on reconnect.

% ===================================================================
\section{Implementation Architecture}
\label{sec:impl}

The reference implementation (Python 3.10+, Pydantic v2, NetworkX) has 23 source modules across four layers: Parsing, Models, Semantics, Runtime.

\textbf{Ontology Registry} (\texttt{adl\_core\_ontology.yaml} v0.2): 5 classes, 10 predicates, 9 actions, 4 scope prefixes. Preconditions use Comparator enum (EQ, NEQ, GT, GTE, LT, LTE, IN, EXISTS) with no \texttt{eval()}.

\textbf{Hybrid Memory}: Hot (in-memory, $<1$ms), Warm (SQLite+NetworkX, 5--20ms), Cold (file-backed). Cascade filtering: status bitmap $\to$ type index $\to$ namespace trie $\to$ vector ANN $\to$ graph traversal.

\textbf{Action Executor}: 6-step pipeline (validate name, check params, evaluate preconditions, execute side effects, trigger transition, update log). Side effects via Protocol plugins.

\textbf{Real-Time Watcher}: Attaches to append()---laundering patterns, status transitions, frequency thresholds, amount alerts.

\textbf{Feishu/Lark Bridge}: Publish, announce, listen, sync-memory, dashboard---enterprise messaging as transport layer.

% ===================================================================
\section{Experimental Evaluation}
\label{sec:eval}

\subsection{Experiment Design}

\begin{table}[h]\centering\caption{Experiments with research question mapping.}\label{tab:exp}
\begin{tabular}{cllc}
\toprule ID & Name & Validation Goal & RQ \\
\midrule
E1 & Chain integrity    & Hash chain + tamper detection  & RQ1 \\
E2 & Status derivation  & Status from chain, not stored  & RQ1 \\
E3 & Snapshot roundtrip & FrontMatter $\leftrightarrow$ EventChain  & RQ1 \\
E4 & Precondition enforce & Comparator pipeline integration & RQ2 \\
E5 & 5-agent audit      & Multi-agent validation + fork   & RQ3 \\
E6 & IBM AML pipeline   & Smart ontology discovery (P=1.000, R=0.500, F1=0.667) & RQ2 \\
\bottomrule
\end{tabular}\end{table}

\subsection{Results}

All six experiments pass. \textbf{E1-E3}: hash chain correctly links events, tampering detected. \textbf{E4}: all 9 Comparator types validated. \textbf{E5}: 5 agents validate with distinct confidence scores. \textbf{E6}: IBM AML HI-Small\_Trans (9,300 transactions), smart heuristic achieves precision 1.000, recall 0.500, F1 0.667.

\subsection{Fork Threshold Sensitivity}

Sweep 0.70--0.95 on 21 controlled concept pairs: \textbf{zero false positives across all thresholds}. Default 0.90 is conservative (F1=0.667, FPR=0.0). For higher recall, 0.75--0.80 achieves F1=0.857.

\subsection{RDF/OWL/PROV-O Mapping}

Turtle compiler: concepts $\to$ named graphs, relations $\to$ owl:ObjectProperty, evidence $\to$ prov:Entity. 100\% information preservation on all 5 example documents (9/9 relations, 10/10 evidence).

\subsection{Baseline Comparison}

8-dimension comparison: ADL provides unique combination of Pydantic validation, SSA pronoun check, cryptographic audit trail (EventChain), ForkManager consensus, YAML registry, L2 prose, Markdown-native authoring---all in a single format.

\subsection{Distributed Deployment (10--100 Agents)}

\begin{table}[h]\centering\caption{Distributed fault scenario results.}\label{tab:dist}
\begin{tabular}{lcccl}
\toprule Scenario & Scale & Writes & Merged & Outcome \\
\midrule
S1: Concurrent & 10/50/100 & 30--300 & 100\% & All events survive \\
S2: Partition  & 3$\times$(10/33/100) & 60--600 & Deterministic & Order-independent \\
S3: Clock skew & 10, $\pm$5hr & 20 & Deterministic & Skew does not affect \\
S4: Adversarial & 1 poison + 9 honest & 68 & Honest confidence 0.80 & Survives attack \\
\bottomrule
\end{tabular}\end{table}

All chains pass \texttt{verify\_integrity()}. Merge is deterministic regardless of entry order. Honest confidence survives adversarial poison events.

% ===================================================================
\section{Discussion}
\label{sec:discussion}

\subsection{Comparison with Object-First Ontologies}

\begin{table}[h]\centering\caption{Object-first vs. event-first.}\label{tab:comparison}
\begin{tabular}{p{3.5cm}p{5cm}p{5cm}}
\toprule Aspect & Object-First (typical) & Event-First (ADL Lite) \\
\midrule
Concept repr.   & Mutable object, stored properties & Append-only event chain \\
State derivation & Explicit writes to fields & Computed from chain \\
Audit trail     & Requires explicit logging & Inherent (cryptographic) \\
Concurrency     & Race conditions & Write-time: conflict-free \\
Trust model     & Implicit in storage & Per-concept hash chain \\
\bottomrule
\end{tabular}\end{table}

\subsection{Applicability to Semantic Web}

ADL Lite bridges LLM authoring and RDF/OWL/PROV-O consumption: concepts $\to$ named graphs, events $\to$ annotations, relation blocks $\to$ OWL assertions.

\subsection{Failure Modes and Adversarial Analysis}
\label{sec:failure}

\textbf{Timestamp collisions}: SHA-256 secondary key for determinism; HLC/vector clocks for production. \textbf{Tampering}: chained hash provides evidence, not prevention; Ed25519 signing planned. \textbf{Chain poisoning}: watcher alerts signal but do not prevent; rate limiting future work. \textbf{SSA circumvention}: creative circumlocutions bypass pattern detection---confirmed experimentally. \textbf{Adversarial validation}: payload tampering detected, 100-event spam triggers alerts, 100-fork bombing handled correctly.

\subsection{CRDT Formalization of Derived States}
\label{sec:crdt}

The derived state forms a join-semilattice $\mathcal{S} = (\text{Status} \times \mathbb{R}_{[0,1]} \times \mathbb{N} \times \mathbb{N}, \sqcup)$:
\[
\text{merge}(s_1, s_2) = \big(\max(\text{status}_1, \text{status}_2),\ \max(\text{conf}_1, \text{conf}_2),\ \max(\text{val}_1, \text{val}_2),\ \max(\text{ev}_1, \text{ev}_2)\big)
\]
Status is lattice-ordered: $\text{PROVISIONAL} \sqsubset \text{FORKED} \sqsubset \text{VALIDATED} \sqsubset \text{DEPRECATED} \sqsubset \text{ARCHIVED}$.

\textbf{Commutativity}: $\text{merge}(A,B) = \text{merge}(B,A)$ --- verified by enumeration. \textbf{Associativity}: $\text{merge}(\text{merge}(A,B),C) = \text{merge}(A,\text{merge}(B,C))$ --- verified over $4^3$ state triples. \textbf{Idempotence}: $\text{merge}(A,A) = A$. \textbf{Monotonicity}: $\text{apply}(s, e) \sqsupseteq s$. \textbf{Convergence}: under partition, $\text{merge}(s_A, s_B) \sqsupseteq s_A$ and $\sqsupseteq s_B$. Verified by 15 executable proofs (\texttt{tests/test\_crdt\_proofs.py}).

\subsection{Limitations}
\label{sec:limitations}

\begin{enumerate}[label=(\arabic*)]
    \item \textbf{LWW semantics (addressed)}: CRDT-correct alternative defined in \S\ref{sec:crdt}.
    \item \textbf{Single-concept granularity}: Cross-concept queries require ADLMemory layer.
    \item \textbf{No formal semantics}: No description logic model theory.
    \item \textbf{Scalability}: Hot layer $\approx$10M skeletons/node; $O(n)$ traversal for long chains.
    \item \textbf{Lark coupling}: Platform-specific; SideEffect protocol enables alternatives.
\end{enumerate}

\subsection{Future Work}

\begin{enumerate}[label=(\arabic*)]
    \item RDF/OWL compiler with named graph mapping.
    \item Distributed EventChain (Kafka-backed).
    \item LLM agent benchmark for ADL authoring.
    \item Robust confidence aggregation (trust-weighted, per-agent credentials).
    \item SHACL integration from ontology registry.
    \item Cross-platform bridges (Slack, Teams, Discord).
    \item Formal CRDT lattice proofs for full state.
    \item Per-agent Ed25519 signatures + Merkle root anchoring.
\end{enumerate}

% ===================================================================
\section{Conclusion}
\label{sec:conclusion}

We presented ADL Lite, an event-first operational ontology for multi-agent KG authoring. The key insight---concepts as append-only event chains---resolves fundamental tensions: race conditions become impossible, audit trails become inherent, trust becomes cryptographic. The four-layer Markdown-native syntax enables LLM agents; SSA constrains ambiguity; the consensus protocol with fork management and offline sync enables autonomous agent fleets. Validated by 6 experiments (6/6), distributed deployment at 10--100 agents, 350 tests at 81\% coverage, CRDT lattice proofs, and a 100\%-accurate RDF/OWL/PROV-O mapping.

\section*{Data Availability}
Source code, ontology registry, experiments, and tests: \url{https://github.com/sunnyang1/adl-lite} (MIT License). Reproduce: \texttt{python -m experiments.runner all} and \texttt{pytest tests/ -v}.

\section*{Acknowledgments}
[TO BE COMPLETED]

\bibliographystyle{unsrt}
\bibliography{references}
\end{document}
